% Canonical TeX edition. Initial text inherited from 3b067cc2eba4ac2f2f91248448727639d8ac7c5a
% Source: research/case_study/phase256_two_stage_rank_four.md
% Subsequent mathematical edits require a new audit.
% Options for packages loaded elsewhere
\PassOptionsToPackage{unicode}{hyperref}
\PassOptionsToPackage{hyphens}{url}
\documentclass[
  11pt,
  a4paper,
]{article}
\usepackage{xcolor}
\usepackage[margin=1in]{geometry}
\usepackage{amsmath,amssymb}
\setcounter{secnumdepth}{-\maxdimen} % remove section numbering
\usepackage{iftex}
\ifPDFTeX
  \usepackage[T1]{fontenc}
  \usepackage[utf8]{inputenc}
  \usepackage{textcomp} % provide euro and other symbols
\else % if luatex or xetex
  \usepackage{unicode-math} % this also loads fontspec
  \defaultfontfeatures{Scale=MatchLowercase}
  \defaultfontfeatures[\rmfamily]{Ligatures=TeX,Scale=1}
\fi
\usepackage{lmodern}
\ifPDFTeX\else
  % xetex/luatex font selection
  \setmainfont[]{STIXGeneral}
  \setmonofont[]{Menlo}
\fi
% Use upquote if available, for straight quotes in verbatim environments
\IfFileExists{upquote.sty}{\usepackage{upquote}}{}
\IfFileExists{microtype.sty}{% use microtype if available
  \usepackage[]{microtype}
  \UseMicrotypeSet[protrusion]{basicmath} % disable protrusion for tt fonts
}{}
\makeatletter
\@ifundefined{KOMAClassName}{% if non-KOMA class
  \IfFileExists{parskip.sty}{%
    \usepackage{parskip}
  }{% else
    \setlength{\parindent}{0pt}
    \setlength{\parskip}{6pt plus 2pt minus 1pt}}
}{% if KOMA class
  \KOMAoptions{parskip=half}}
\makeatother
\setlength{\emergencystretch}{3em} % prevent overfull lines
\providecommand{\tightlist}{%
  \setlength{\itemsep}{0pt}\setlength{\parskip}{0pt}}
\usepackage{newunicodechar}
\newfontfamily\numeralfallback{STIX Two Text}
\newunicodechar{∎}{\ensuremath{\blacksquare}}
\usepackage{fvextra}
\DefineVerbatimEnvironment{Highlighting}{Verbatim}{breaklines,breakanywhere,breakanywheresymbolpre={},breakanywheresymbolpost={},breaksymbolleft={},breaksymbolright={},commandchars=\\\{\}}
\newunicodechar{⁰}{{\numeralfallback\char"2070\relax}}
\newunicodechar{₀}{{\numeralfallback\char"2080\relax}}
\newunicodechar{₁}{{\numeralfallback\char"2081\relax}}
\newunicodechar{₂}{{\numeralfallback\char"2082\relax}}
\newunicodechar{₃}{{\numeralfallback\char"2083\relax}}
\newunicodechar{₄}{{\numeralfallback\char"2084\relax}}

\usepackage{seqsplit}
\renewcommand{\texttt}[1]{{\ttfamily\seqsplit{#1}}}
\setlength{\emergencystretch}{6em}

\newunicodechar{⁻}{{\numeralfallback\char"207B\relax}}
\newunicodechar{⁸}{{\numeralfallback\char"2078\relax}}
\newunicodechar{ⁱ}{{\numeralfallback\char"2071\relax}}
\newunicodechar{⁴}{{\numeralfallback\char"2074\relax}}
\newunicodechar{⁷}{{\numeralfallback\char"2077\relax}}
\newunicodechar{⁶}{{\numeralfallback\char"2076\relax}}
\newunicodechar{⁵}{{\numeralfallback\char"2075\relax}}
\newunicodechar{₈}{{\numeralfallback\char"2088\relax}}
\newunicodechar{ᵢ}{\textsubscript{i}}
\newunicodechar{ⱼ}{\textsubscript{j}}
\newunicodechar{ₜ}{\textsubscript{t}}
\usepackage{bookmark}
\IfFileExists{xurl.sty}{\usepackage{xurl}}{} % add URL line breaks if available
\urlstyle{same}
% fallback for those not using the hyperref driver hyperxmp:
\makeatletter
\@ifundefined{xmpquote}{\newcommand{\xmpquote}[1]{#1}}{}
\makeatother
\hypersetup{
  hidelinks,
  pdfcreator={LaTeX via pandoc}}

\author{}
\date{}

\begin{document}

\section{Two homology covers in rank
four}\label{two-homology-covers-in-rank-four}

Author manuscript, phase256, 11 October 2026. This candidate proof is
prepared for independent review and is not yet admitted. Its conclusion
concerns finite groups of rank four. It does not answer the finite-cover
question in higher rank. The earlier arbitrary-cover theorem in phase216
is separate.

\subsection{1. Statement and
dependencies}\label{statement-and-dependencies}

A marked string group of rank four is a group generated by distinct
nonidentity involutions s\_0,s\_1,s\_2,s\_3, with s\_i and s\_j
commuting whenever \textbar i−j\textbar\textgreater1. Write
G\_I=⟨s\_i:i∈I⟩. Throughout this manuscript G is finite and

\begin{verbatim}
s_0 ∉ G123,                 s_3 ∉ G012.                 (1)
\end{verbatim}

No intersection property is assumed. Put

\begin{verbatim}
g=|G|, A=G012, B=G123, C=G12, q=ord(s_1s_2).
\end{verbatim}

The distinctness assumptions imply q≥2 and
\textbar G02\textbar=\textbar G03\textbar=\textbar G13\textbar=4.

\textbf{Theorem 1 (candidate).} For every prime p\textgreater4g there is
a finite marked string C-group T mapping onto G whose kernel is a
p-group.

A finite polytope's monodromy group satisfies (1): the subgroup
generated by all marks except s\_i preserves the rank-i face of every
flag, whereas s\_i changes that face. Thus Theorem 1 would imply the
following conclusion.

\textbf{Corollary 2 (candidate).} Every finite rank-four polytope whose
monodromy group is not a string C-group has infinitely many pairwise
nonisomorphic finite minimal regular polytopal covers. For every prime
p\textgreater4\textbar Mon(P)\textbar, one can choose such a cover whose
degree over Mon(P) is a nontrivial power of p.

Problem 5 in Cunningham--Pellicer, \emph{Open problems on k-orbit
polytopes}, §3.1, asks whether non-C monodromy can coexist with a unique
minimal regular polytopal cover. The question has no finiteness
restriction. Cunningham, \emph{Chiral polytopes whose smallest regular
cover is a polytope}, §4, expresses the same expectation and a more
specific family conjecture. Corollary 2 asserts a finite-cover
strengthening for finite polytopes of rank four; it is not a claim to
settle the finite strengthening in all ranks.

The construction has two steps. First take the mod-p homology cover of a
surface with boundary; this repairs the two rank-three ends. Then fill
the middle boundary polygons and take another mod-p homology cover. The
new argument proves the precise complement conditions needed in this
second step, including the nonorientable case.

The proof reuses the cellular intersection argument and closed-surface
homology-cover theorem of phase211, §§3--5. Those arguments are restated
at the point where they are used. The raw surface and its relative
collapses were developed in the unfinished
phase224\_finite\_rank\_four.md, §§2--5; this manuscript supplies the
previously missing uniform second step. The known existence of finite
regular polytopal covers is due to Monson--Schulte. Mere existence is
not the new conclusion asserted here. The new assertion is the existence
of a cover with a p-group kernel for every sufficiently large prime, and
its consequence for finite minimal covers.

\subsection{2. The surface with boundary and the first
cover}\label{the-surface-with-boundary-and-the-first-cover}

Let W be generated by r\_0,r\_1,r\_2,r\_3, with defining relations
r\_i²=1 and {[}r\_0,r\_2{]}={[}r\_0,r\_3{]}={[}r\_1,r\_3{]}=1. Let
K=ker(W→G).

Construct a polygonal complex X with vertex set G and a colored edge
from x to xs\_i for every i. Attach one square for each left coset of
G02, G03, and G13 along its alternating two-color cycle. The squares
have embedded boundaries. Each edge of color 0 or 3 lies in two squares;
each edge of color 1 or 2 lies in one. The link of every vertex is the
path

\begin{verbatim}
2 — 0 — 3 — 1.
\end{verbatim}

Consequently X is a connected compact surface with boundary. Its
boundary circles are exactly the alternating 12 cycles, one for each
left coset of C. Its cell counts give

\begin{verbatim}
χ(X)=g−2g+3g/4=−g/4.
\end{verbatim}

The path presentation gives π\_1(X,1)=K: backtracking gives the
involution relations, and square homotopies give all conjugates of the
three commuting relations. One square per coset suffices, since its
boundary based at another corner is a conjugate of that boundary or its
inverse.

For I⊆\{0,1,2,3\}, let X\_I be the connected colored subcomplex through
1, including exactly the available squares whose colors belong to I. Its
vertex set is G\_I. Set X\_A=X012 and X\_B=X123.

Fix a prime p and put

\begin{verbatim}
V=H_1(X;F_p), Φ=[K,K]K^p, H=W/Φ.
\end{verbatim}

The subgroup Φ is characteristic in K and hence normal in W. The surface
cover X\_p→X with deck group V is connected and has the Cayley graph of
H as its one-skeleton. Since a compact connected surface with nonempty
boundary has H\_2=0 over every field,

\begin{verbatim}
dim V=1+g/4,              |H|=g p^(1+g/4).             (2)
\end{verbatim}

All marks and (1) lift to H. Write A\_H=H012, B\_H=H123, and C\_H=H12.

\textbf{Lemma 3.} The groups A\_H and B\_H, with their indicated marks,
are string C-groups.

\textbf{Proof.} For S=X01∪X12, every 02 square has a 2-edge outside S:
otherwise s\_0∈G12. Collapse that square across such an edge; each
2-edge belongs to only one square. Every 13 square has a 1-edge outside
S: all vertices of S belong to A, so two such edges in S would imply
s\_3∈A. Collapse these squares too. Each remaining 03 square has a
0-edge outside S, for the same reason. The 02 square formerly incident
with that edge was already removed, so it is now a free edge and the 03
square can be collapsed across it. These collapses fix S and leave only
a graph attached to S. Therefore H\_1(S;F\_p)→H\_1(X;F\_p) is injective.
The dual argument applies to X12∪X23.

Here is the homological intersection argument used for these pairs. For
connected subcomplexes E,F with a common base point, suppose H\_1(E∩F)=0
and H\_1(E∪F)→H\_1(X) is injective. If paths in E and F lift from the
chosen base point to the same vertex, their difference is zero in
H\_1(X), hence in H\_1(E∪F). The Mayer--Vietoris connecting map first
shows that their common projected endpoint belongs to the base component
C\_0 of E∩F. Comparing both paths with a path in C\_0 and using
H\_1(E∩F)=0 then shows that they have the same lifted endpoint as that
path. Thus the chosen lifted components intersect exactly in the chosen
lift of C\_0.

For the two pairs above, E∩F is a disjoint union of single edges, so
this argument gives H01∩H12=H1 and H12∩H23=H2. A group generated by two
distinct involutions has the rank-two intersection property. The
recursive criterion for string C-groups now proves the lemma. ∎

\subsection{3. Removing an end and detecting boundary
classes}\label{removing-an-end-and-detecting-boundary-classes}

We first record two elementary connectivity facts with their cell
details.

\textbf{Lemma 4.} The induced Cayley graph on G\textbackslash A is
connected. The same holds for G\textbackslash B, and for the
corresponding ends in any finite marked group satisfying (1).

\textbf{Proof.} Contract each left A-coset, retaining the edges of color
3 between cosets. The resulting finite graph is connected and
vertex-transitive under left multiplication by G. It has at least two
vertices. A finite connected vertex-transitive graph has no cut vertex:
a finite connected graph with at least two vertices has a noncut vertex
(for example, a leaf of a spanning tree), and transitivity carries this
property to every vertex. Deleting the vertex A therefore leaves a
connected graph. Every remaining A-coset is internally connected using
colors 0,1,2, so a path in the contracted graph lifts to a path avoiding
A. The dual argument proves the other assertion. Loops or repeated edges
in the contracted graph do not affect this argument. ∎

\textbf{Lemma 5.} The following complements have the same components as
their induced Cayley graphs on vertices outside the indicated
subcomplexes:

\begin{enumerate}
\def\labelenumi{(\alph{enumi})}
\item
  X\textbackslash X\_A and X\textbackslash X\_B;
\item
  Y(Y012∪Y123), where Y is obtained from X by filling every middle
  boundary circle with a disk.
\end{enumerate}

Here the analogous statement holds for any finite marked group
satisfying (1).

\textbf{Proof.} An outside face always has an outside vertex, and the
outside vertices of that face are connected along its boundary through
outside vertices. We check this assertion explicitly. For X\_A, a 02
square is wholly in X\_A or has no vertex in A. A 03 square meets A in
one 0-edge or not at all; a 13 square meets it in one 1-edge or not at
all. These are all possibilities because s\_3∉A. Thus the asserted
boundary connectivity holds. The argument for X\_B is dual.

For the union of the two ends, a 02 square lying outside Y012 has no A
vertex; it meets B in one 2-edge or not at all because s\_0∉B. A 13
square is dual. A 03 square meets A in at most one 0-row and B in at
most one 3-column; removing such a row and column leaves either the
whole boundary, a connected edge, or a single vertex. A middle 12 disk
is wholly in the union or has no vertex there, since C belongs to both
ends.

We also need that every outside edge has an outside endpoint. For the
single end A, edges of colors 0,1,2 with an A endpoint are selected,
whereas an edge of color 3 cannot have both endpoints in A. For the top
union, an edge of color 1 or 2 with a selected endpoint is selected. An
outside edge of color 0 has no endpoint in A and cannot have both
endpoints in B, because s\_0∉B. Color 3 is dual. Thus the required
assertion holds.

At a vertex outside the selected subcomplex every incident edge and face
is outside it. The link of that vertex is connected, including the
boundary case in X. Hence all incident outside faces belong to the same
component. Conversely, whenever outside faces meet across an outside
edge, or around an outside vertex, their outside vertices are connected
by the boundary paths just described. Thus the components of the outside
face adjacency graph, the topological complement, and the induced
outside-vertex graph are in bijection. ∎

For a boundary circle xC, choose either orientation and denote its class
in V by u\_x. To view this as the class of a based loop one may use any
path from 1 to x; the path and its inverse cancel in first homology. The
choice of orientation only changes the sign. Define

\begin{verbatim}
V_A=image(H_1(X_A;F_p)→V),
V_B=image(H_1(X_B;F_p)→V).
\end{verbatim}

\textbf{Lemma 6.} For every x∈G,

\begin{verbatim}
u_x∈V_A ⇔ x∈A,             u_x∈V_B ⇔ x∈B.           (3)
\end{verbatim}

\textbf{Proof.} If x∈A, the boundary circle belongs to X\_A and is
joined to 1 inside X\_A, so u\_x∈V\_A. Suppose x∉A. By Lemmas 4--5 the
complement of X\_A is connected, including its outside face adjacency
graph. The number of middle boundary circles outside A is

\begin{verbatim}
[G:C]−[A:C]=([G:A]−1)[A:C]≥2,                     (4)
\end{verbatim}

because A is proper and s\_0∉C implies {[}A:C{]}≥2.

If u\_x∈V\_A, an oriented boundary cycle b\_x differs from a cycle z
supported in X\_A by the boundary of a cellular two-chain c over
F\_p.~Across every edge outside X\_A which is not a boundary edge of X,
the coefficients of the incident outside squares must agree up to the
incidence sign, since b\_x−z has coefficient zero there. Choose an edge
on an outside boundary circle different from xC; its coefficient in
b\_x−z is zero and it is incident with just one square. The coefficient
of that square in c must therefore be zero. Connectivity of the outside
face adjacency graph propagates zero to every outside square. But an
edge of b\_x also has just one incident square, which is outside X\_A,
and its coefficient in b\_x−z is nonzero. This is a contradiction. The
argument is valid over every field; it makes no orientability
assumption. The proof for B is dual. ∎

\subsection{4. Avoidance of the chosen end
components}\label{avoidance-of-the-chosen-end-components}

Choose the lift of 1 corresponding to the identity of H. The chosen
lifts of X\_A and X\_B have vertex sets A\_H and B\_H. Above a vertex
x∈A, the chosen A-component meets the fiber in one affine coset of V\_A;
above x∉A it has no vertex. Indeed, the differences of two lifted
endpoints of A-paths are exactly the homology classes of loops in X\_A.
The analogous assertion holds for B and V\_B. Choosing an origin in each
fiber identifies the deck action with translation by V; all these
assertions are independent of those origins.

Let E be the induced graph on H(A\_H∪B\_H). We call a vertex whose image
lies in G(A∪B) an exterior-base vertex. The distinction matters: E also
contains vertices over A∪B which do not lie in the two chosen
components.

\textbf{Lemma 7.} Assume p\textgreater4g. Then E is connected.

\textbf{Proof.} A union of two proper subgroups cannot equal a group: if
neither contains the other, the product of an element in
A\textbackslash B with one in B\textbackslash A lies in neither. Thus
there is a vertex h∈G(A∪B).

At any exterior-base vertex x, the whole middle circle xC avoids A∪B.
All its lifts therefore lie in E. Traversing it changes the deck
coordinate by u\_x, which belongs to neither V\_A nor V\_B by Lemma 6.
Consequently one can translate any point over x by k u\_x, for every
k∈F\_p, using paths in E.

Consider a fixed base path γ of length L from x, and fix an initial
lift. Translate its initial point by k u\_x; the entire lifted path is
translated by that same deck vector. For any visited vertex over A, the
condition that this translated vertex lies in the chosen A-component is
an affine equation k u\_x∈a+V\_A. Since u\_x∉V\_A, it has at most one
solution k. A visited vertex over B excludes at most one further value.
Thus at most 2(L+1) values of k are forbidden. Whenever
p\textgreater2(L+1), one can first travel around the safe middle circle
to choose a permitted translate and then lift γ wholly in E.

Every exterior-base vertex reaches some point over h: take a base path
to h of length at most g−1 and use the preceding argument. We next show
that the entire fiber over h is connected in E. Choose a spanning tree
of the base Cayley graph rooted at h. Its fundamental graph loops have
length at most 2g−1, and their classes generate H\_1(X;F\_p), since
attaching squares only takes a quotient of graph homology. Let γ be any
one of these loops, with class v, and let z be any point over h. There
is a permitted k because p\textgreater4g≥2(length(γ)+1). Travel from z
by k u\_h around the middle circle, lift γ, and then travel by −k u\_h
around the middle circle at h. All three paths lie in E, and their net
deck displacement is v. Thus every spanning-tree loop class can be
realized between points of the fiber by paths in E. Those classes
generate V, so the whole fiber is connected.

It remains to reach an exterior-base vertex from an arbitrary vertex z
of E. Let x be its image. If x∈A\textbackslash B, the edge of color 3
takes x to a vertex outside both A and B: s\_3∉A, while s\_3∈B and x∉B.
The lifted target is therefore automatically allowed. The case
x∈B\textbackslash A is dual, using color 0. If x∈A∩B, first use color 0.
The lift remains outside the chosen A-component because it stays in its
current A-component, which is unchosen. Its new base vertex xs\_0 lies
in A\textbackslash B because s\_0∉B, so it is outside the chosen
B-component too. Apply the preceding color-3 step. This proves that
every vertex of E reaches the already connected exterior-base region. ∎

The field bound is deliberately coarse. No enumeration of H, whose order
is given in (2), is required by this argument.

\subsection{5. Orientability and the two local
complements}\label{orientability-and-the-two-local-complements}

Let Y\_H be the closed surface obtained by filling every middle boundary
circle in X\_p.~Equivalently, it is constructed directly from H by
adding 02,03,13 and12 polygons. It satisfies all the hypotheses of the
surface construction in phase211. Lemmas 5 and 7 show that

\begin{verbatim}
Y_H\((Y_H)012∪(Y_H)123)                           (5)
\end{verbatim}

is connected. We must still treat its orientability and the two smaller
selected subcomplexes; connectedness alone is insufficient in odd
characteristic for a nonorientable ambient surface.

\textbf{Lemma 8.} If Y\_H is nonorientable, the complement in (5) is
nonorientable.

\textbf{Proof.} For both the raw and capped surface, orientation
transport along a colored edge reverses the local orientation convention
determined by the common colored link. Hence the orientation character
of a closed graph walk is its length modulo two. In particular X is
orientable exactly when its Cayley graph is bipartite. Capping boundary
circles does not change orientability, since every middle boundary has
even length.

If X is orientable, its cover X\_p and its capped surface Y\_H are
orientable, so only the other case needs consideration. Suppose X is
nonorientable. For a spanning tree rooted at h as in Lemma 7, at least
one fundamental graph loop γ has odd length; otherwise every graph cycle
would have even length. Its length is at most 2g−1. The construction in
Lemma 7 realizes its class v between any point z over h and z+v by a
path in E. The two added middle walks have even length, so the
constructed path has odd length. Repeat the construction p times, making
a new permitted choice of k at each stage. The total displacement is
pv=0 in V. Since p is odd, this gives a closed odd-length walk in E. Its
orientation character is nontrivial, so E lies in a nonorientable
component of the complement. That complement is connected by (5), hence
is nonorientable. This argument also shows directly that Y\_H is
nonorientable in this case. ∎

\textbf{Lemma 9.} Each of

\begin{verbatim}
Y_H\((Y_H)01∪(Y_H)12),
Y_H\((Y_H)12∪(Y_H)23)                            (6)
\end{verbatim}

is connected. If Y\_H is nonorientable, both are nonorientable.

\textbf{Proof.} We treat the first. Its selected vertex set S\_v=H01∪H12
lies in A\_H. By Lemma 4, the induced graph on H\textbackslash A\_H is
connected. Every vertex of A\_H\textbackslash S\_v has a color-3
neighbor outside A\_H, so the induced graph on H\textbackslash S\_v is
connected too.

We verify that its components are the components of the surface
complement. For a 02 square, H01 meets its vertices in a 0-edge or not
at all, whereas H12 meets them in a 2-edge or not at all. Here Lemma 3
ensures s\_2∉H01, and extreme independence ensures s\_0∉H12. Removing
one such row and column leaves connected outside vertices. A 03 square
meets H01 in at most a 0-edge and H12 in at most one vertex. For the
latter assertion, neither s\_0, s\_3 nor s\_0s\_3 belongs to H12: the
first two are excluded by (1), and s\_0s\_3∈H12⊆B\_H would imply
s\_0∈B\_H. Again the remaining vertices are nonempty and connected on
the boundary. A 13 square meets each subgroup in a 1-edge or not at all.
Their selected edges cannot be opposite rows, because all selected
vertices lie in A\_H and s\_3∉A\_H.

Finally a 12 polygon is wholly selected if it is the chosen H12 coset.
Otherwise its intersection with H01, when nonempty, is a coset of
H01∩H12=H1 by Lemma 3, hence a single boundary edge of color 1. Deleting
its two vertices leaves a connected boundary path, since the polygon has
at least four vertices. Thus every outside face has outside vertices
connected along its boundary. Every outside edge also has an outside
endpoint: a color-1 edge with a selected endpoint is selected; an
outside color-0 edge has no H01 endpoint and cannot have both endpoints
in H12; an outside color-2 edge has no H12 endpoint and cannot have both
endpoints in H01, by the end intersection property; and a color-3 edge
cannot have both endpoints in A\_H. The connected-link argument of Lemma
5 therefore completes the comparison and proves connectivity. The other
complement is dual.

The top complement (5) is contained in each complement (6). If Y\_H is
nonorientable, Lemma 8 supplies an orientation-reversing loop in (5),
hence in each of the connected complements (6). They are therefore
nonorientable. ∎

\subsection{6. The second cover and the minimal
descendants}\label{the-second-cover-and-the-minimal-descendants}

For completeness, the surface criterion being used is the following. If
Y is a connected closed surface and S a connected proper subcomplex, the
map H\_1(S;F\_p)→H\_1(Y;F\_p) is injective in either of these
situations:

\begin{itemize}
\tightlist
\item
  Y is orientable and Y\textbackslash S is connected;
\item
  p is odd, Y is nonorientable, and every component of Y\textbackslash S
  is nonorientable.
\end{itemize}

Indeed the exact sequence of the pair reduces the assertion to
surjectivity of H\_2(Y)→H\_2(Y,S). Coefficients on outside faces
propagate across outside edges. In the orientable case, one coefficient
is free on each component, and the ambient fundamental class supplies
their common diagonal. In odd characteristic a nonorientable component
forces its coefficient to equal its negative and therefore to vanish.
These observations give the stated criterion, including precisely the
needed coefficient qualification.

Apply it to Y\_H and all three selected subcomplexes in (5)--(6). Lemmas
7--9 show that it applies. In each of the three intersections of the
paired subcomplexes, every component is a disk or a single edge: the top
pair has one filled 12 disk for each common middle coset, while the
local pairs have edges of color 1 or 2. Their first homology is zero.
The homological intersection argument from Lemma 3 therefore applies to
the mod-p homology cover of Y\_H and gives all three equalities

\begin{verbatim}
T012∩T123=T12,   T01∩T12=T1,   T12∩T23=T2.
\end{verbatim}

The recursive string C-group criterion proves that T is a string
C-group. The path presentation of Y\_H identifies its fundamental group
with the kernel of W(q\_H)→H, where q\_H is the order of the marked
middle product in H and W(q\_H)=W/⟨⟨(r\_1r\_2)\^{}\{q\_H\}⟩⟩. The kernel
of its mod-p abelianization is characteristic, so the resulting deck
cover is the Cayley complex of a finite marked group T mapping onto H,
with elementary abelian p-kernel. Its composite T→H→G consequently has a
p-group kernel. This proves Theorem 1, subject to independent review of
the complete argument.

Now suppose G is the non-C monodromy group of a finite rank-four
polytope P. The standard string C-group/polytope correspondence makes T
a finite regular polytopal cover of P. Explicitly, write the flags of P
as W/J, where J is a base-flag stabilizer, so K=core\_W(J). A regular
polytope with normal flag kernel N covers P if and only if N⊆J,
equivalently N⊆K. Normality makes this condition independent of the
chosen base flag. Thus every regular quotient still covering P has its
marked group mapping onto G=W/K. Among its regular quotients still
covering P, choose one minimal under covering; finiteness of T ensures
its existence. Call its marked group Q\_p.~The kernel of Q\_p→G is a
quotient of the p-group ker(T→G), hence is a p-group. It is nontrivial,
since otherwise G itself would be a string C-group. Thus
\textbar Q\_p\textbar/\textbar G\textbar{} is a nontrivial power of p.
Different primes give different group orders, and therefore different
regular polytopes. There are infinitely many primes above 4g, proving
Corollary 2.

The same argument also separates the universal kernels: for distinct
primes p and ℓ, the common quotient over G of the two constructed covers
has kernel whose order divides both a p-power and an ℓ-power, so that
kernel is trivial. Their kernel product in W is therefore exactly K.
This last observation is not needed for the distinct-order proof of the
corollary.

\subsection{7. Attribution, verification, and current
limits}\label{attribution-verification-and-current-limits}

The surface intersection criterion, the careful treatment of the hidden
common deck kernel, and the passage from p-group covers to finite
minimal descendants reuse phase211 and the unfinished raw-surface
calculations in phase224\_finite\_rank\_four.md. The newly asserted
steps are the boundary membership test (Lemma 6), the large-prime path
construction (Lemma 7), and its combination with orientability and the
local complements (Lemmas 8--9). Standard graph connectivity,
covering-space theory, cellular homology, Mayer--Vietoris, and the
string C-group correspondence are used as stated.

The exact program phase256\_voltage\_check.py checks five explicitly
supplied marked groups of orders 32, 96, 96, 216 and 120 at primes 131,
389, 389, 877 and 487, respectively. It computes the raw homology
quotient, tests (3) at every base vertex, and checks translated paths
for three specified deck coordinates and every fundamental graph loop.
Four of these groups fail the original one-stage closed-surface
criterion. The order-120 example is A5×C2 and is nonorientable; the
program also constructs a closed odd-length exterior walk by repeating
the permitted loop 487 times. These checks use symbolic deck coordinates
and do not enumerate the enormous first covers. The four other inputs
are marked-group diagnostics, with no asserted realization as monodromy
groups of actual polytopes.

The Lean file HomologyCoverAvoidance.lean proves five elementary
statements: a transverse affine line meets one affine coset at most
once; fewer affine constraints than field elements leave a permitted
scalar; the specialization to two possible constraints per visited
vertex; the path-length inequality; and preservation of odd parity under
odd repetition. It uses only Lean's standard logical axioms. The replay
also rejects removal of transversality, weakening of the strict field
bound, and replacement of odd repetition by even repetition. These
formal statements verify the algebra and counting step used in Lemmas
7--8. Cellular topology, the identification of the actual end fibers
with affine cosets, and the final polytope correspondence remain written
proofs. Neither the finite checks nor these five formal interfaces alone
proves Theorem 1.

The unrestricted finite-cover question in rank at least five remains
open within this work. This manuscript makes no assertion about new
classifications of polytopes or about how much earlier literature it
replaces.

\end{document}
